../cictikz/src/cictikz/data/tex/cictikz_preamble.tex

% The switch symbols: an NMOS pass device, a PMOS one, the transmission gate
% built from both, and the bowtie symbol that stands for it.
%
% The original draws the single devices as simplified three terminal switches
% where only the arrow direction separates N from P. Real circuitikz symbols
% are used instead, so the PMOS carries its bubble and the two cannot be
% confused. The control labels are the original's: c on the NMOS and c-bar on
% the PMOS, so both conduct when c is high.
%
% On the bowtie the inversion is carried by the bubble on the upper control
% line rather than by a bar on the label, which is how the original draws it.

\definecolor{axisblue}{rgb}{0.16,0.16,0.55}
\definecolor{portgreen}{rgb}{0.00,0.45,0.20}

% A labelled terminal: position, anchor, label.
\newcommand{\swPort}[3]{%
  \draw[axisblue, thick] #1 circle (0.09);
  \node[portgreen, anchor=#2] at #1 {#3};
}

\begin{circuitikz}[american, transform shape]

  % ---- NMOS pass device ----
  \def\xa{0}
  \draw[axisblue, thick] (\xa+0.1,0) to [Tnmos, n=MN] (\xa+2.0,0);
  \draw[axisblue, thick] (MN.gate) -- ++(0,0.55);
  \swPort{(\xa+0.0,0)}{east}{$A$\;}
  \swPort{(\xa+2.1,0)}{west}{\;$B$}
  \swPort{($(MN.gate)+(0,0.55)$)}{south}{$c$}

  % ---- PMOS pass device, complementary control ----
  \def\xb{4.0}
  \draw[axisblue, thick] (\xb+0.1,0) to [Tpmos, n=MP] (\xb+2.0,0);
  \draw[axisblue, thick] (MP.gate) -- ++(0,0.55);
  \swPort{(\xb+0.0,0)}{east}{$A$\;}
  \swPort{(\xb+2.1,0)}{west}{\;$B$}
  \swPort{($(MP.gate)+(0,0.55)$)}{south}{$\overline{c}$}

  % ---- Transmission gate, both devices in parallel ----
  \def\xc{8.4}
  \draw[axisblue, thick] (\xc+0.1,0) -- (\xc+0.6,0) -- (\xc+0.6,0.9);
  \draw[axisblue, thick] (\xc+0.6,0.9) to [Tpmos, n=TP] (\xc+2.4,0.9);
  \draw[axisblue, thick] (\xc+2.4,0.9) -- (\xc+2.4,0) -- (\xc+2.9,0);
  \draw[axisblue, thick] (\xc+0.6,0) -- (\xc+0.6,-0.9);
  \draw[axisblue, thick] (\xc+0.6,-0.9) to [Tnmos, n=TN, mirror] (\xc+2.4,-0.9);
  \draw[axisblue, thick] (\xc+2.4,-0.9) -- (\xc+2.4,0);
  \fill[axisblue] (\xc+0.6,0) circle (0.07);
  \fill[axisblue] (\xc+2.4,0) circle (0.07);
  \draw[axisblue, thick] (TP.gate) -- ++(0,0.5);
  \draw[axisblue, thick] (TN.gate) -- ++(0,-0.5);
  \swPort{(\xc+0.0,0)}{east}{$A$\;}
  \swPort{(\xc+3.0,0)}{west}{\;$B$}
  \swPort{($(TP.gate)+(0,0.5)$)}{south}{$\overline{c}$}
  \swPort{($(TN.gate)+(0,-0.5)$)}{north}{$c$}

  % ---- The symbol that stands for all of that ----
  % Inversion lives on the bubble, not on the label, as in the original.
  \def\xd{9.0}
  % Clear of the transmission gate's lower control port above it.
  \def\yd{-5.0}
  \draw[axisblue, thick] (\xd,\yd+0.75) -- (\xd,\yd-0.75) -- (\xd+1.05,\yd) -- cycle;
  \draw[axisblue, thick] (\xd+2.1,\yd+0.75) -- (\xd+2.1,\yd-0.75) -- (\xd+1.05,\yd) -- cycle;
  \draw[axisblue, thick] (\xd-0.6,\yd) -- (\xd,\yd);
  \draw[axisblue, thick] (\xd+2.1,\yd) -- (\xd+2.7,\yd);
  \draw[axisblue, thick] (\xd+1.05,\yd+1.05) -- (\xd+1.05,\yd+0.40);
  \draw[axisblue, thick] (\xd+1.05,\yd+0.28) circle (0.12);
  \draw[axisblue, thick] (\xd+1.05,\yd) -- (\xd+1.05,\yd-1.05);
  \swPort{(\xd-0.7,\yd)}{east}{$A$\;}
  \swPort{(\xd+2.8,\yd)}{west}{\;$B$}
  \swPort{(\xd+1.05,\yd+1.15)}{south}{$c$}
  \swPort{(\xd+1.05,\yd-1.15)}{north}{$c$}

\end{circuitikz}
\end{document}
